% PLANTILLA DE ACTIVIDAD — INVESTIGACIÓN APLICADA A LA CONTADURÍA / UAS
\documentclass[spanish,letterpaper,oneside]{article}
\def\documenttitle {Planeación de actividad}
\def\documentsubtitle {Investigación aplicada a la contaduría}
\def\documentsubject {Actividad [POR COMPLETAR]}
\def\documentauthor {Nombre del estudiante}
\def\coursename {Investigación aplicada a la contaduría}
\def\coursecode {Clave por confirmar}
\def\universityname {Universidad Autónoma de Sinaloa}
\def\universityfaculty {Licenciatura en Contaduría}
\def\universitydepartment {Unidad académica por confirmar}
\def\universitydepartmentimage {UAS/assets/marca-uas-provisional}
\def\universitydepartmentimagecfg {height=1.45cm}
\def\universitylocation {Sinaloa, México}
\def\authortable {
  \begin{tabular}{ll}
    \textbf{Estudiante:} & Nombre por completar \\
    Matrícula: & Por completar \\
    \textit{Docente:} & Nombre por completar \\
    Grupo y periodo: & Por completar \\
    Fecha: & \today
  \end{tabular}
}
\input{template}
\usepackage{helvet}
\renewcommand{\familydefault}{\sfdefault}
\setcitestyle{authoryear,open={(},close={)}}
\definecolor{uasBlue}{HTML}{123B6D}
\newcommand{\porcompletar}[1]{\textcolor{uasBlue}{\textbf{[POR COMPLETAR: #1]}}}

\begin{document}
\templatePortrait
\templatePagecfg
\onehalfspacing
\begin{abstractd}
  Ficha editable para traducir una consigna de aula en un producto verificable, congruente y trazable.
\end{abstractd}
\templateIndex
\templateFinalcfg

\section{Identificación de la actividad}
\begin{description}
  \item[Producto solicitado:] \porcompletar{reporte, protocolo, matriz, instrumento, exposición u otro.}
  \item[Consigna original:] \porcompletar{archivar en la carpeta de planeaciones y resumir sin alterar requisitos.}
  \item[Fecha y formato:] \porcompletar{fecha, extensión, nomenclatura y medio de entrega.}
  \item[Fuentes obligatorias:] \porcompletar{registrar únicamente fuentes reales en el archivo bibliográfico.}
\end{description}

\section{Propósito y problema contable}
\porcompletar{explicar qué debe demostrar el producto y qué problema de la contaduría activa.}

\section{Entregables y evidencias}
\begin{longtable}{@{}p{0.25\linewidth}p{0.42\linewidth}p{0.25\linewidth}@{}}
\toprule
\textbf{Entregable} & \textbf{Contenido verificable} & \textbf{Criterio de aceptación} \\
\midrule
\endfirsthead
\toprule
\textbf{Entregable} & \textbf{Contenido verificable} & \textbf{Criterio de aceptación} \\
\midrule
\endhead
\bottomrule
\endlastfoot
\porcompletar{producto} & \porcompletar{evidencia requerida} & \porcompletar{rúbrica o condición} \\
\end{longtable}

\section{Ruta de trabajo}
\begin{enumerate}
  \item Delimitar la pregunta y el objetivo de la actividad.
  \item Localizar fuentes verificables y registrar sus datos bibliográficos.
  \item Seleccionar método, técnica o estructura conforme al producto solicitado.
  \item Construir la evidencia y distinguir datos, análisis y postura.
  \item Revisar la congruencia y los criterios de la rúbrica.
\end{enumerate}

\section{Criterios de autoevaluación}
\begin{itemize}
  \item Cumplimiento íntegro de la consigna sin copiarla como contenido.
  \item Pertinencia para la práctica contable y claridad del problema.
  \item Uso responsable de datos, confidencialidad y fuentes.
  \item Correspondencia entre objetivos, método, evidencia y conclusión.
  \item Redacción, presentación y citación consistentes.
\end{itemize}

\section{Notas y pendientes}
\porcompletar{registrar dudas para el docente, datos faltantes y decisiones tomadas.}

\bibliography{investigacion-aplicada-a-la-contaduria}
\end{document}
